\begin{abstract}
  %
  Transformers are the driving force behind today's Large Language
  Models (LLMs), serving as the foundation for their performance and
  versatility.
  %
  Yet, their compute and memory costs grow with sequence length,
  posing scalability challenges for long-context inferencing.
  %
  In response, the algorithm community is exploring alternative
  architectures—such as state space models (SSMs) (e.g., Mamba-2),
  linear attention, and recurrent neural networks (RNNs)—which we
  refer to as \emph{post-transformers}.
  %
  This shift presents a key challenge: building a serving system that
  efficiently supports not only emerging post-transformer LLMs but
  also existing transformer models within a unified framework.
  %

  %
  To address this challenge, we analyze the performance
  characteristics of transformer and post-transformer LLMs.
  %
  Despite their algorithmic differences, both are largely
  bounded by memory bandwidth under batched inference—due to
  attention in transformers and state updates in post-transformers.
  %
  Inspired by this finding, we propose \sysname{}, an accelerator
  solution that aims to address the memory bottleneck by jointly
  leveraging (1) Processing-in-Memory (PIM) paradigm and (2) LLM quantization.
  %
  Further analyses suggest two additional insights: (1) state update
  operations, unlike attention, incur high hardware cost, making
  per-bank PIM acceleration inefficient, and (2) different
  low-precision arithmetic methods offer varying accuracy-area
  tradeoffs, while we identify Microsoft's MX as a Pareto-optimal choice.
  %
  Building on these insights, we design the architecture of
  \sysname{} as an array of
  State-update Processing Units (SPUs), each shared between two banks
  to enable interleaved access.
  %
  Each SPU includes a State-update Processing Engine (SPE) that
  comprises element-wise multipliers and adders using MX-based
  quantized arithmetic, enabling efficient execution of state update
  and attention operations.
  %
  Our evaluation shows that, compared to LLM-optimized GPU and
  GPU+PIM systems, \sysname{} achieves up to 4.1$\times$ and
  2.1$\times$ higher generation throughput, respectively.
  %

\end{abstract}
